% 本片：第 4 篇（Emerton--Gee--Hellmann）印刷页 359--370（PDF 页 372--383）。
% 上片 ch-a4-16 覆盖印刷页 347--358（PDF 页 360--371）；其 p.371 末行为无号公式 $\chi=\bigotimes_v\chi_{\delta_v,\mathrm{sm}}:A^+(G)\to L^\times$，
% 本片自 p.372 首行句中断续处起笔，故首句承接该公式之后的「of the Atkin--Lehner ring.」。

\ysrc{359}

……Atkin--Lehner 环的。此外，我们有特征空间的如下等同
\begin{equation}\tag{9.6.39}\label{eq:a4-17-9-6-39}
M^\dagger(G,K^p,(\xi_v)_{v\mid p})[\frakm_\rho]^{A^+(G)=\chi}\cong\operatorname{Hom}_G\Bigl(\widehat{\otimes}_v\mathcal F^G_B(\delta_v),\hat S(K^p,L)^{\mathrm{an}}[\frakm_\rho]\Bigr).
\end{equation}

设 $\xi_v$ 是与正则 Hodge--Tate 权 $\underline\lambda_v$ 相伴的代数特征标。则局部解析表示 $\mathrm{Ind}^{I_v}_{B_{v,0}}z^{\xi_v}$ 包含最高权为 $\xi_v$ 的代数表示 $V_{\underline\lambda_v}$。特别地，权为 $\xi=(z^{\xi_v})_{v\mid p}$ 的超收敛 $p$-进自守形式空间 $M^\dagger(G,K^p,(z^{\xi_v})_{v\mid p})$ 包含由权为 $\xi$ 的经典自守形式组成的空间
\[
S(K^pI,(\kappa_v)_v)=\left\{\,f:G(F)\backslash G(\mathbf A_F^{\infty})/K^p\to\bigotimes_{v\mid p}V_{\underline\lambda_v}\ \middle|\
\begin{array}{l}
\text{满足}\ f(gu)=u^{-1}f(g)\\[2pt]
\text{对一切}\ g\in G(\mathbf A_F^{\infty}),\ u\in\prod_{v\mid p}I_v
\end{array}\,\right\}.
\]

\begin{lem}{9.6.40}\label{lem:a4-17-9-6-40}
在等同 (9.6.39) 下，一个超收敛 $p$-进自守形式
\[
f\in M^\dagger(G,K^p,(z^{\xi_v})_{v\mid p})[\frakm_\rho]^{A^+(G)=\chi}
\]
是经典的（即属于 $S(K^pI,(z^{\xi_v})_v)$），当且仅当相应的态射
\[
\widehat{\otimes}_v\mathcal F^G_B(\delta_v)\longrightarrow\hat S(K^p,L)^{\mathrm{an}}[\frakm_\rho]
\]
经过局部代数商\footnote{关于该商映射的描述见注 6.2.21。}
\[
\widehat{\otimes}_v\mathcal F^G_B(\delta_v)\longrightarrow\bigotimes_{v\mid p}V_{\underline\lambda_v}\otimes\bigl(\mathrm{Ind}^{G(F_v)}_{B_v}\delta_{\mathrm{sm}}\bigr)^{\mathrm{sm}}
\]
分解。
\end{lem}

\begin{proof}
其证明与 [BHS17a, Proposition 3.4] 的证明类似。
\end{proof}

下述猜想基本上由标准的 Fontaine--Mazur 型猜想推出。

\begin{conj}{9.6.41}\label{conj:a4-17-9-6-41}
设 $\rho\in\XX_{G,F,S}$ 是具正则 Hodge--Tate 权 $\underline\lambda$ 的晶体表示\tn{底本此处在花体 $\mathfrak X$ 与 $\mathcal X$ 之间混用（本页印作 $\mathfrak X_{G,F,S}$）；本译本全篇统一记作 $\XX$（即 $\mathcal X$），见下方 $\ol{\XX}_B$ 等处。}，$\chi$ 是 Atkin--Lehner 环的特征标，使得特征空间
\[
M^\dagger(G,K^p,(z^{\xi_v})_{v\mid p})[\frakm_\rho]^{A^+(G)=\chi}
\]
非零，即 $\rho$ 与某个 $p$-进自守形式 $f$ 相伴。则特征空间
\[
S(K^pI,(z^{\xi_v})_v)[\frakm_\rho]^{A^+(G)=\chi}
\]
非零，即 $\rho$ 与某个经典自守形式相伴。
\end{conj}

鉴于我们关于「从局部解析表示到伽罗瓦表示叠上的凝聚层」的函子所提出的猜想，人们可以通过研究相应的凝聚层态射来逼近这一猜想：把猜想中的函子 $\mathfrak A^{\mathrm{rig}}_G$ 作用于局部解析表示之间的态射
\[
\widehat{\otimes}_v\mathcal F^G_B(\delta_v)\longrightarrow\bigotimes_{v\mid p}\bigl(V_{\underline\lambda_v}\otimes\mathrm{Ind}^{G(F_v)}_{B_v}\delta_{\mathrm{sm}}\bigr)^{\mathrm{sm}},
\]
得到凝聚层之间的一个态射，再研究该态射。事实上，在定理 9.6.36 那样的额外 Taylor--Wiles 假设下，这一策略已经给出猜想 9.6.41 的证明，见 [BHS19, Theorem 5.1.3]。

\ysrc{360}

猜想 9.6.41 中的 $p$-进自守形式 $f$ 是否自动是经典的？这是一个有趣的问题。注意伴点或伴形式的现象并不涵盖这一点：伴形式是非经典权的 $p$-进形式（因而绝不可能是经典形式！），而这里我们已经把权固定成经典权。尽管如此，我们的猜想仍蕴含：不能指望 $f$ 是经典的。猜想 9.6.18 与注 9.6.25 用可由 Drinfeld 紧化 $\ol{\XX}_B$ 计算的凝聚层，给出如下特征空间的精确猜想性描述：
\[
S(K^pI,(z^{\xi_v})_v)[\frakm_\rho]^{A^+(G)=\chi}\qquad\text{与}\qquad M^\dagger(G,K^p,(z^{\xi_v})_{v\mid p})[\frakm_\rho]^{A^+(G)=\chi}.
\]
结果表明，该紧化一般不是 Gorenstein 的；这些奇点带来的后果是层 $\mathcal M^*(G,K^p)$ 的纤维维数可能跳跃。这基本上意味着
\[
M^\dagger(G,K^p,(z^{\xi_v})_{v\mid p})[\frakm_\rho]^{A^+(G)=\chi}
\]
有时应当具有比其经典子空间更大的维数。在三个变元的定酉群情形，这一预期可以做成具体例子，见 [HHS25]。

\section*{附录 A.\quad 范畴论背景}

我们回顾 $(\infty)$-范畴论中若干将在正文中使用的结论。

\subsection*{A.1.\quad 稳定 $\infty$-范畴}

在范畴局部 Langlands 纲领中，我们所考虑的函子必然是导出的：它们未必把 $p$-进群 $G$ 的光滑表示阿贝尔范畴中的对象，映成 Langlands 参数模叠上凝聚层或拟凝聚层阿贝尔范畴中的对象；相反，某些表示会映成层的复形。（与 Kisin--Taylor--Wiles 拼接的相容性要求已经迫使我们如此，因此这是整个故事中不可避免的一部分。）出于这一原因，我们必须使用适当的层复形范畴。

如果我们在二十世纪写这些讲义，接下来就会说明：这让我们不得不使用各种各样的导出范畴或三角范畴，然后逐一说明是哪些。但我们将利用同调代数在当代的进展，改用稳定 $\infty$-范畴。这类范畴避开了三角范畴众所周知的技术缺陷，并且当我们想施行各种粘接过程或取极限的过程时尤其好用（这两种过程我们都会用到：例如我们处理的是形式代数叠上的层，而形式代数叠既有拓扑结构——需要粘接论证——又有 Ind-叠结构，需要取极限的论证）。

下面我们极简要地回顾其中一些基础，主要目的是为那些对传统导出范畴与三角范畴理论比对稳定 $\infty$-范畴理论更熟悉的读者（以及作者！）指个方向。为处理集合论问题，我们遵循 [Lur09a] 的进路。特别地，我们固定一个 Grothendieck 宇宙；属于这个固定宇宙的集合称为小的。此外，一切极限与余极限都假定是小极限与小余极限，下文一般不再说明。

\ysrc{361}

稳定 $\infty$-范畴及其上 $t$-结构的基础，参看 [Lur17, Ch.~1]。若 $\mathcal C$ 是稳定 $\infty$-范畴，则其底层的同伦范畴带有典范的三角范畴结构。实践中往往可以直接设想自己在底层的三角范畴中工作，只是有一个改进：锥（从而同伦极限与同伦余极限）是典范的。

一个重要的例子是 Grothendieck 阿贝尔范畴 $\mathcal A$ 所对应的导出 $\infty$-范畴 $D(\mathcal A)$，它的同伦范畴就是 $\mathcal A$ 的通常（无界）导出范畴。（回忆：称 $\mathcal A$ 为 Grothendieck 阿贝尔范畴，是指它余完备、$\mathcal A$ 中滤过余极限的构造正合，并且 $\mathcal A$ 有一组生成元。）本文通篇采用上同调编号，而 [Lur17]（以及我们引用的大多数文献）采用同调编号。相应地，我们讨论上链复形，而所引用的结论则关于链复形；对此一般不再说明。记住这一点：导出 $\infty$-范畴 $D(\mathcal A)$ 在 [Lur17, Def.~1.3.5.8] 中定义为如下全子范畴的微分分次神经——该全子范畴由 $\mathcal A$ 中上链复形范畴里对某个模型结构而言为纤维性的对象组成。（该模型结构的构造见 [Lur17, Prop.~1.3.5.3]，纤维性对象的部分描述见 [Lur17, Prop.~1.3.5.6]。）我们也可以把 $D(\mathcal A)$ 描述为：上链复形范畴关于拟同构的 $\infty$-范畴局部化。（见 [Lur17, Prop.~1.3.5.15] 及其证明。）

若 $X$ 与 $Y$ 是 $\infty$-范畴 $\mathcal C$ 的对象，则有映射 anima\footnote{[Lur09b] 与 [Lur17] 用「空间」（spaces）这一术语指单纯集在通常 Quillen 模型结构下产生的 $\infty$-范畴及其对象；我们沿用 [CS24] 的术语 anima。} $\mathrm{Maps}_{\mathcal C}(X,Y)$。当 $\mathcal C$ 还是稳定的时，这个映射 anima 实际上是连通谱。在我们的语境中，范畴还会是 $\mathcal O$-线性的（$\mathcal O$ 是某个系数环），因此 $\mathrm{Maps}_{\mathcal C}(X,Y)$ 是连通的 $\mathcal O$-模谱；等价地说，它是一个动态（animated）$\mathcal O$-模，由 Dold--Kan 对应可以看作 $D^{\leq0}(\mathcal O\text{-}\mathrm{Mod})$ 的对象，而 $D^{\leq0}(\mathcal O\text{-}\mathrm{Mod})$ 是以 $\mathcal O$-模的非正分次上链复形之导出范畴为底层的稳定 $\infty$-范畴。

再次利用稳定性，还可以定义一个（通常非连通的）谱 $\mathrm{RHom}_{\mathcal C}(X,Y)$，由它可通过公式
\[
\mathrm{Maps}_{\mathcal C}(X,Y)=\tau^{\leq0}\mathrm{RHom}_{\mathcal C}(X,Y)
\]
恢复出映射空间。同样，在我们考虑的语境中范畴是 $\mathcal O$-线性的，因此各个 $\mathrm{RHom}_{\mathcal C}(X,Y)$ 都是 $\mathcal O$-模谱，可以把它们看作 $D(\mathcal O\text{-}\mathrm{Mod})$ 的对象（两个方向都无界）。

在 $\mathcal C$ 上放一个 $t$-结构（按定义）相当于在其底层的三角范畴上放一个 $t$-结构。于是，对 $n\in\Z$，全子范畴\footnote{当然，指的是子 $\infty$-范畴！} $\mathcal C^{\leq n}$、$\mathcal C^{\geq n}$ 定义为同伦范畴中相应子范畴的拉回。关键事实如下：
\begin{enumerate}[label=(\arabic*)]
\item 若 $X\in\mathcal C^{\leq n}$ 且 $Y\in\mathcal C^{\geq n+1}$，则 $\mathrm{Maps}_{\mathcal C}(X,Y)=0$（即可缩）；
\item 含入 $\mathcal C^{\leq n}\hookrightarrow\mathcal C$ 有右伴随 $\tau^{\leq n}$；
\item 含入 $\mathcal C^{\geq n}\hookrightarrow\mathcal C$ 有左伴随 $\tau^{\geq n}$。
\end{enumerate}

\ysrc{362}

\begin{defn}{A.1.1}\label{defn:a4-17-A-1-1}
若 $\mathcal C$ 是赋有 $t$-结构的稳定 $\infty$-范畴，则定义 $\mathcal C$ 的中心（heart）为交
\[
\mathcal C^{\heartsuit}:=\mathcal C^{\geq0}\cap\mathcal C^{\leq0}.
\]
\end{defn}

中心 $\mathcal C^{\heartsuit}$ 是阿贝尔范畴 [Lur17, Rem.~1.2.1.12]。（更确切地说，它等价于 $\mathcal C$ 的同伦范畴上相应 $t$-结构的中心之神经，而后者按字面就是阿贝尔范畴。回忆：在 [Lur09a] 的框架中，取神经是把通常范畴转成 $\infty$-范畴的技术手段。）

若 $\mathcal A$ 是 Grothendieck 阿贝尔范畴，则稳定 $\infty$-范畴 $D(\mathcal A)$ 带有典范的 $t$-结构，它来自 $\mathcal A$ 的导出范畴上的通常 $t$-结构 [Lur17, Prop.~1.3.5.21]。存在全忠实函子 $\mathcal A\hookrightarrow D(\mathcal A)$，它诱导 $\mathcal A$ 与 $D(\mathcal A)$ 上 $t$-结构之中心之间的等价。（按定义，$D(\mathcal A)$ 的中心与其全子范畴 $D^+(\mathcal A)$（见下文）的中心一致，于是该断言可由 [Lur17, Prop.~1.3.2.19] 的对偶版本推出。）我们有时把 $\mathcal A$ 称为 $D(\mathcal A)$ 中由静态对象组成的子范畴（静态对象有时也叫「离散」对象或「经典」对象）。

可以构造全子范畴 $D^+(\mathcal A)=\bigcup_{n\in\Z}D^{\geq n}(\mathcal A)$。这个范畴也可以构造为：$\mathcal A$ 的内射对象的有下界上链复形所成全子范畴的微分分次神经。（这是因为有下界的纤维性复形恰是有下界的内射复形；见 [Lur17, Prop.~1.3.5.6]。于是 $D^+(\mathcal A)$ 的这一构造，与 [Lur17, §1.3.2] 中对有足够多投射对象的阿贝尔范畴给出的 $D^-(\mathcal A)$ 的构造相对偶。另注意：$D^+(\mathcal A)$ 的这一构造只要求 $\mathcal A$ 有足够多内射对象，而不要求 $\mathcal A$ 是 Grothendieck 的。）

\subsection*{A.2.\quad 结合代数对象}

$E_1$-环（有时称为 $E_1$-环谱或 $A_\infty$-环谱）的概念在 [Lur17, Defn.~7.1.0.1] 中定义。由 [Lur17, Rem.~7.1.0.3]，静态的 $E_1$-环就是通常意义的结合环。关于 $E_1$-环 $A$ 上的左模与右模有一套理论，我们分别用 $\mathrm{LMod}_A$、$\mathrm{RMod}_A$ 记这些稳定 $\infty$-范畴。当 $A$ 是静态的时，它们等同于通常模阿贝尔范畴的导出 $\infty$-范畴，[Lur17, Rem.~7.1.1.16]。

若 $x$ 是稳定 $\infty$-范畴 $\mathcal C$ 的对象，则由 [Lur17, Rem.~7.1.2.2]，谱 $\mathrm{RHom}_{\mathcal C}(x,x)$ 获得 $E_1$-环的结构，我们记之为 $\End_{\mathcal C}(x)$；特别地，对一切 $n\in\Z$ 有 $\pi_n\End_{\mathcal C}(x)=\Ext^{-n}_{\mathcal C}(x,x)$。

\subsection*{A.3.\quad Pro-范畴（与 Ind-范畴）}

我们极简要地回顾 Pro- 与 Ind-$\infty$-范畴的定义；详见 [Lur09a, §5.3.5]。（更确切地说，该文献处理的是 Ind-$\infty$-范畴的情形，但对任意 $\infty$-范畴 $\mathcal C$，有等价 $\Ind(\mathcal C)\cong\mathrm{Pro}(\mathcal C^{\mathrm{op}})^{\mathrm{op}}$。）

若 $\mathcal C$ 是（小）$\infty$-范畴，可以定义与之相伴的 pro-范畴 $\mathrm{Pro}(\mathcal C)$ 如下：其对象是以余滤过小 $\infty$-范畴 $I$ 为指标范畴的图 $F:I\to\mathcal C$；两个图 $F:I\to\mathcal C$ 与 $G:J\to\mathcal C$ 之间的态射由公式
\[
\operatorname{Mor}_{\mathrm{Pro}(\mathcal C)}(F,G)=\varprojlim_{J}\varinjlim_{I}\operatorname{Mor}_{\mathcal C}(F(i),G(j))
\]
定义。（事实上，正如 [Lur09a, §5.3] 的引言所解释的，这并不是建立 Ind- 与 Pro-$\infty$-范畴理论最方便的方式，但我们忽略这一点。）

\ysrc{363}

当 $\mathcal C$ 是通常范畴时，我们有时用「$\varprojlim_I$」$F(i)$ 记图 $F:I\to\mathcal C$，并称之为 $\mathcal C$ 的一个 pro-对象。这一记号的意义在于：把该 pro-对象与 $\mathcal C$ 中的极限 $\varprojlim_I F(i)$（若该极限碰巧存在）区分开来。采用这一记号时，我们常写 $X_i$ 而不写 $F(i)$，从而把 $\mathrm{Pro}(\mathcal C)$ 的对象 $F$ 记作「$\varprojlim_I$」$X_i$。

若 $\mathcal C$ 是稳定 $\infty$-范畴，则 $\Ind\mathcal C$ 与 $\mathrm{Pro}\,\mathcal C$ 也是。若 $\mathcal C$ 带有 $t$-结构，则 $\Ind\mathcal C$ 与 $\mathrm{Pro}\,\mathcal C$ 上有典范诱导的 $t$-结构（例如见 [AGH19, Prop.~2.13]）。

\subsection*{A.4.\quad 连续函子}

\begin{defn}{A.4.1}\label{defn:a4-17-A-4-1}
设 $\mathcal C$ 是容许滤过余极限的 $\infty$-范畴。若 $\mathrm{Maps}_{\mathcal C}(x,-)$ 保持滤过余极限，则称对象 $x\in\mathcal C$ 是紧的（compact）（见 [Lur09a, Defn.~5.3.4.5]）。
\end{defn}

\begin{rem}{A.4.2}\label{rem:a4-17-A-4-2}
若 $\mathcal C$ 是容许滤过余极限的稳定 $\infty$-范畴，则它余完备。此时 $x\in\mathcal C$ 紧，当且仅当 $\mathrm{RHom}_{\mathcal C}(x,-)$ 保持滤过余极限，当且仅当 $\mathrm{RHom}_{\mathcal C}(x,-)$ 保持一切（小）余极限。（这是因为 $\mathrm{RHom}_{\mathcal C}(x,-)$ 自动保持一切极限，从而是正合的。）
\end{rem}

\begin{defn}{A.4.3}\label{defn:a4-17-A-4-3}
若 $\infty$-范畴 $\mathcal C$ 容许滤过余极限、其紧对象子范畴 $\mathcal C^c$ 是小范畴，且自然函子 $\Ind(\mathcal C^c)\to\mathcal C$（把 $\mathcal C^c$ 中的滤过图送到它在 $\mathcal C$ 中的余极限）是等价，则称 $\mathcal C$ 是紧生成的（compactly generated）。
\end{defn}

\begin{defn}{A.4.4}\label{defn:a4-17-A-4-4}
若稳定 $\infty$-范畴之间的函子保持（小）余极限，则称它是连续的。
\end{defn}

\begin{rem}{A.4.5}\label{rem:a4-17-A-4-5}
在许多文献中，「连续」（continuous）指保持滤过余极限；因此我们称为连续的稳定 $\infty$-范畴之间的函子，通常会被叫作既正合又连续。由于我们没有机会考虑非正合函子，我们选择把正合性条件并入连续函子的定义。
\end{rem}

笔记 [Cam16] 对与紧对象相关的范畴论概念给出了有用的概述。下述引理部分基于这些笔记中讨论的一些结果。

\begin{lem}{A.4.6}\label{lem:a4-17-A-4-6}
设 $F:\mathcal C\to\mathcal D$ 是紧生成稳定 $\infty$-范畴之间的函子。
\begin{enumerate}[label=(\arabic*)]
\item $F$ 连续，当且仅当它有右伴随 $G$。
\item 若 $F$ 连续，则 $F$ 保持紧对象，当且仅当其右伴随 $G$ 连续。
\item 若 $F$ 连续、全忠实且保持紧对象，则 $F$ 也反射紧对象，即 $F(X)$ 紧蕴含 $X$ 本身紧。
\end{enumerate}
\end{lem}

\begin{proof}
第一部分是伴随函子定理 [Lur09a, Cor.~5.5.2.9]。对第 \textup{(2)} 部分，先设 $F$ 保持紧对象。为证 $G$ 连续，只需证它保持有限余极限与滤过余极限。对前者，回忆稳定 $\infty$-范畴之间的函子保持有限余极限，当且仅当它保持有限极限；而 $G$ 是右伴随，事实上保持一切极限。

\ysrc{364}

于是对 $\mathcal D$ 中任意滤过余极限 $\varinjlim_j Y_j\xrightarrow{\sim}Y$，我们须证诱导态射 $\varinjlim_j G(Y_j)\to G(Y)$ 是同构。由 Yoneda 引理，只需对 $\mathcal C$ 的任意对象 $X$ 作用 $\operatorname{Mor}(X,-)$ 之后验证；又因 $\mathcal C$ 紧生成，可进一步假定 $X$ 紧。此时有
\begin{align*}
\operatorname{Mor}\bigl(X,\varinjlim_j G(Y_j)\bigr)
&=\varinjlim_j\operatorname{Mor}\bigl(X,G(Y_j)\bigr)
\xrightarrow{\sim}\varinjlim_j\operatorname{Mor}\bigl(F(X),Y_j\bigr)\\
&\xrightarrow{\sim}\operatorname{Mor}\bigl(F(X),\varinjlim_j Y_j\bigr)
\xrightarrow{\sim}\operatorname{Mor}\bigl(F(X),Y\bigr)
\xrightarrow{\sim}\operatorname{Mor}\bigl(X,G(Y)\bigr),
\end{align*}
即所求。

反之，若 $G$ 连续，则它特别保持滤过余极限。于是若 $X$ 紧、$\varinjlim_j Y_j$ 是 $\mathcal D$ 中的滤过余极限，则可写
\begin{align*}
\operatorname{Mor}\bigl(F(X),\varinjlim_j Y_j\bigr)
&\xrightarrow{\sim}\operatorname{Mor}\bigl(X,G(\varinjlim_j Y_j)\bigr)
\xrightarrow{\sim}\operatorname{Mor}\bigl(X,\varinjlim_j G(Y_j)\bigr)\\
&=\varinjlim_j\operatorname{Mor}\bigl(X,G(Y_j)\bigr)
\xrightarrow{\sim}\varinjlim_j\operatorname{Mor}\bigl(F(X),Y_j\bigr),
\end{align*}
故 $F(X)$ 紧，即所求。

最后看 \textup{(3)}。注意到若 $F(X)$ 紧，则有
\begin{align*}
\operatorname{Mor}\bigl(X,\varinjlim_j Y_j\bigr)
&\xrightarrow{\sim}\operatorname{Mor}\bigl(F(X),F(\varinjlim_j Y_j)\bigr)
\xrightarrow{\sim}\operatorname{Mor}\bigl(F(X),\varinjlim_j F(Y_j)\bigr)\\
&=\varinjlim_j\operatorname{Mor}\bigl(F(X),F(Y_j)\bigr)
\xrightarrow{\sim}\varinjlim_j\operatorname{Mor}\bigl(X,Y_j\bigr),
\end{align*}
故 $X$ 紧，即所求。
\end{proof}

\subsection*{A.5.\quad 完备 $t$-结构}

本小节的参考文献是 [BZNP17, §6] 与 [Lur17, §§1.2, 1.3]。（重要提醒：这些文献使用同调编号，而我们用的是上同调编号！）

\begin{defn}{A.5.1}\label{defn:a4-17-A-5-1}
若 $\mathcal C$ 是赋有 $t$-结构的稳定 $\infty$-范畴，则它有左完备化
\[
\wh{\mathcal C}:=\varprojlim\mathcal C^{\geq n}.
\]
该极限以全序集 $\Z$ 为指标，过渡函子由截断 $\tau^{\geq n}:\mathcal C^{\geq n-1}\to\mathcal C^{\geq n}$ 给出。

存在典范函子 $\mathcal C\to\wh{\mathcal C}$，它把对象 $X$ 映到序列 $(\tau^{\geq n}X)$，后者自然地可以解释为 $\wh{\mathcal C}$ 的对象。若该函子是等价，则称 $\mathcal C$（连同其给定的 $t$-结构）是左完备的。

对偶地有右完备化与右完备性的概念。我们不引入右完备化的记号，但明确写出：若由 $X\mapsto(\tau^{\leq n}X)$ 定义的典范函子 $\mathcal C\to\varprojlim\mathcal C^{\leq n}$ 是等价，则称 $\mathcal C$ 是右完备的。

与上述完备性概念相关的是相应的分离性概念。
\end{defn}

\begin{defn}{A.5.2}\label{defn:a4-17-A-5-2}
若 $\mathcal C$ 是赋有 $t$-结构的稳定 $\infty$-范畴，且
\[
\mathcal C^{-\infty}:=\bigcap_{n\in\Z}\mathcal C^{\leq n}
\]
只含零对象，则称它是左分离的。

\ysrc{365}

对偶地，若
\[
\mathcal C^{\infty}:=\bigcap_{n\in\Z}\mathcal C^{\geq n}
\]
只含零对象，则称它是右分离的。
\end{defn}

\begin{rem}{A.5.3}\label{rem:a4-17-A-5-3}
若带有 $t$-结构的稳定 $\infty$-范畴 $\mathcal C$ 左完备，则 $\mathcal C$ 左分离。（若 $X$ 是 $\mathcal C^{-\infty}$ 的对象，则对一切 $n$ 有 $\tau^{\geq n}X=0$，于是在等价 $\mathcal C\xrightarrow{\sim}\wh{\mathcal C}$ 下，对象 $X$ 等同于一个零对象序列，即 $\wh{\mathcal C}$ 的零对象；故 $X$ 是 $\mathcal C$ 的零对象。）
\end{rem}

若 $\mathcal A$ 是 Grothendieck 阿贝尔范畴，则稳定 $\infty$-范畴 $D(\mathcal A)$ 既左分离又右分离，并且右完备 [Lur17, Prop.~1.3.5.21]。它一般不必左完备 [Nee11]。

若 $\mathcal A$ 容许可数积，并且这些积还是正合的，则 $D(\mathcal A)$ 容许可数积（在复形上逐层取积），且由 [Lur17, Prop.~1.2.1.19] 可知 $D(\mathcal A)$ 左完备（因为该处条件 (2) 满足）。例如，若 $\mathcal A$ 是环上模范畴，则 $D(\mathcal A)$ 既左完备又右完备。我们将关心这样的语境：$\mathcal A$ 中的积虽然存在，却未必正合。因此我们记下 [Lur17, Prop.~1.2.1.19] 的下述温和推广。（这种推广的可能性见 [BZNP17, Rmk.~6.1.5]，那里只要求积的右 $t$-正合性「在有限平移意义下」成立。本质上相同的结果也作为 [Ant21, Prop.~8.14] 得到证明。）

\begin{prop}{A.5.4}\label{prop:a4-17-A-5-4}
设 $\mathcal C$ 是带有 $t$-结构的稳定 $\infty$-范畴。假设 $\mathcal C$ 容许可数积，且可数积的构造具有有界振幅（即存在 $a\geq0$，使得若 $\{X_n\}$ 是 $\mathcal C^{\leq0}$ 中任意对象序列，则 $\prod_n X_n$ 落在 $\mathcal C^{\leq a}$ 中）。则 $\mathcal C$ 左完备，当且仅当 $\mathcal C$ 左分离。
\end{prop}

\begin{proof}
沿用 [Lur17, Prop.~1.2.1.19] 的证明，只是现在注意：若 $\cdots\to X_n\to X_{n-1}\to\cdots$ 是 $\mathcal C^{\leq0}$ 中对象构成的塔，则 $\lim X_n$ 落在 $\mathcal C^{\leq a+1}$ 中。于是该证明照原样可行，只是我们考虑的不是分解 $\lim f\to f(n-1)\to f(n)$，而是分解 $\lim f\to f(n-a-1)\to f(n)$。
\end{proof}

\begin{rem}{A.5.5}\label{rem:a4-17-A-5-5}
我们对上面提到的 [Nee11] 的构造稍作展开：它给出 Grothendieck 阿贝尔范畴 $\mathcal A$ 的例子，使得命题 A.5.4 的假设（$D(\mathcal A)$ 中的可数积具有有界上同调振幅）与结论（$D(\mathcal A)$ 左完备）都失败。下文的例 E.2.3 给出该构造的一个显式例子，它自然地出现在（非解析的）pro-$p$ 群的光滑表示论语境中。
\end{rem}

与三角范畴的框架相比，稳定 $\infty$-范畴框架的优点之一，是下述关于右导出函子存在性的简单而非常一般的结论（它是 [Lur17, Thm.~1.3.3.2] 的对偶版本）。回忆：设 $F:\mathcal C\to\mathcal D$ 是带有 $t$-结构的稳定 $\infty$-范畴之间的函子。若 $F(\mathcal C^{\geq0})\subseteq\mathcal D^{\geq0}$，则称 $F$ 左 $t$-正合；若 $F(\mathcal C^{\leq0})\subseteq\mathcal D^{\leq0}$，则称 $F$ 右 $t$-正合；若二者兼有，则称 $F$ 为 $t$-正合。

\ysrc{366}

\begin{thm}{A.5.6}\label{thm:a4-17-A-5-6}
设 $\mathcal A$ 是有足够多内射对象的阿贝尔范畴，$\mathcal C$ 是带有右完备 $t$-结构的稳定 $\infty$-范畴。则 $F\mapsto\tau^{\leq0}F|_{\mathcal A}$（其中的限制通过把 $\mathcal A$ 与 $D^+(\mathcal A)$ 的中心等同来取）是如下两个范畴之间的等价：一边是「把 $\mathcal A$ 的内射对象送入 $\mathcal C^{\heartsuit}$ 的左 $t$-正合函子 $D^+(\mathcal A)\to\mathcal C$」所成的 $\infty$-范畴，另一边是从 $\mathcal A$ 到 $\mathcal C^{\heartsuit}$ 的左正合函子所成的通常范畴。
\end{thm}

作为一个特别推论，我们有下述结果 [Lur17, Prop.~1.3.3.7]（众所周知，在三角范畴的语境中，由于锥的构造不具有函子性，该结果在同样的一般性程度下不成立）。

\begin{prop}{A.5.7}\label{prop:a4-17-A-5-7}
假设 $\mathcal C$ 是带有右完备 $t$-结构的稳定 $\infty$-范畴，且其中心 $\mathcal C^{\heartsuit}$ 有足够多内射对象。则含入 $\mathcal C^{\heartsuit}\hookrightarrow\mathcal C$ 本质唯一地延拓为一个 $t$-正合函子 $F:D^+(\mathcal C^{\heartsuit})\to\mathcal C$\footnote{意思是：函子 $\infty$-范畴 $\operatorname{Fun}(D^+(\mathcal C^{\heartsuit}),\mathcal C)$ 中由这类延拓组成的子 $\infty$-范畴是可缩的；其 1-态射由那些在 $F|_{\mathcal C^{\heartsuit}}$ 上限制为恒等的自然变换组成，其更高态射由这些 1-态射之间的任意更高态射组成。这就是「在 $F|_{\mathcal C^{\heartsuit}}$ 上限制为恒等的前提下、在同构意义下唯一」这一通常说法的 $\infty$-范畴类比。}。此外，下述性质等价：
\begin{enumerate}[label=(\arabic*)]
\item 该函子全忠实。
\item 对 $\mathcal C^{\heartsuit}$ 中任意对象 $X$ 与 $Y$，自然映射 $\Ext^i_{\mathcal C^{\heartsuit}}(X,Y)\to\Ext^i_{\mathcal C}(X,Y)$ 对 $i>0$ 是同构。
\item 对 $\mathcal C^{\heartsuit}$ 中任意对象 $X$ 与 $Y$，若 $Y$ 内射，则对 $i>0$ 有 $\Ext^i_{\mathcal C}(X,Y)=0$。
\item 对 $\mathcal C^{\heartsuit}$ 中任意对象 $X$ 与 $Y$，若 $Y$ 内射，则存在满态射 $Z\to X$，使得对 $i>0$ 有 $\Ext^i_{\mathcal C}(Z,Y)=0$。
\end{enumerate}
最后，若这些等价性质成立，则 $F$ 的本质像等于 $\mathcal C^+:=\bigcup_n\mathcal C^{\geq n}$。
\end{prop}

\begin{proof}
除插入第 \textup{(2)} 条之外（我们加入它只是为了表述更清楚），这不过是 [Lur17, Prop.~1.3.3.7] 的对偶版本。为看出 \textup{(2)} 与其他条件的等价性，注意 \textup{(1)} 当然蕴含 \textup{(2)}（即使 $i=0$ 也如此；不过这一情形无论如何是自动的，因为按构造 $\mathcal C^{\heartsuit}$ 在 $\mathcal C$ 中是全子范畴），而 \textup{(2)} 显然蕴含 \textup{(3)}。
\end{proof}

下文我们将既按上述形式使用这一命题，也使用它对投射对象与上有界范畴的对偶形式（[Lur17, Prop.~1.3.3.7] 中字面处理的就是这一形式）。

\subsection*{A.6.\quad 紧对象与凝聚对象、正则 $t$-结构}

本小节的参考文献是 [BZNP17, §6]。（提醒：该文献使用同调编号，而非我们用的上同调编号。）

下述内容取自 [BZNP17, Def.~6.2.2]。

\begin{defn}{A.6.1}\label{defn:a4-17-A-6-1}
设 $\mathcal C$ 是容许滤过余极限的稳定 $\infty$-范畴，并带有与滤过余极限相容的 $t$-结构（也就是说：$\mathcal C^{\geq0}$，等价地每个 $\mathcal C^{\geq n}$，在 $\mathcal C$ 中对滤过余极限的构造封闭）。称对象 $x\in\mathcal C$ 是凝聚的（coherent），是指 $x$ 有下界（即对某个 $n$，$x$ 落在 $\mathcal C^{\geq n}$ 中），并且对一切 $m\leq n$，$x$ 在 $\mathcal C^{\geq m}$ 中紧。
\end{defn}

\ysrc{367}

\begin{defn}{A.6.2}\label{defn:a4-17-A-6-2}
在定义 A.6.1 的语境下，我们用 $\Coh(\mathcal C)$ 记 $\mathcal C$ 的凝聚对象所成全子 $\infty$-范畴。由 [BZNP17, Rem.~6.2.3] 它是稳定的子 $\infty$-范畴。
\end{defn}

\begin{rem}{A.6.3}\label{rem:a4-17-A-6-3}
在定义 A.6.1 的语境下，若 $m\geq n$，则从 $x$ 到 $\mathcal C^{\geq m}$ 中对象的映射都经过 $\tau^{\geq m}x$ 分解；因此若 $x$ 在 $\mathcal C^{\geq n}$ 中紧，便知对 $m\geq n$，$\tau^{\geq m}x$ 作为 $\mathcal C^{\geq m}$ 的对象是紧的。
\end{rem}

\begin{rem}{A.6.4}\label{rem:a4-17-A-6-4}
若 $\mathcal C$ 如定义 A.6.1，并且还右完备，又若 $x\in\mathcal C$ 凝聚，则 $x$ 既有下界也有上界。（把紧性假设用于恒等映射 $x\to x=\varinjlim\tau^{\leq n}x$ 即可。）
\end{rem}

下述结果综合了 [BZNP17, Lems.~6.2.4, 6.2.5] 的部分内容。

\begin{lem}{A.6.5}\label{lem:a4-17-A-6-5}
设 $\mathcal C$ 是带有右完备 $t$-结构、且该 $t$-结构与滤过余极限相容的稳定 $\infty$-范畴。则下述条件等价：
\begin{enumerate}[label=(\arabic*)]
\item 含入 $\mathcal C^{\geq0}\hookrightarrow\mathcal C^{\geq-1}$ 保持紧对象。
\item $\mathcal C^{\heartsuit}$（即 $\mathcal C$ 上 $t$-结构的中心）的紧对象所成全子范畴是 $\mathcal C^{\heartsuit}$ 的阿贝尔子范畴。
\end{enumerate}
若这些条件成立，则 $\mathcal C$ 中的凝聚对象恰是那些既有上界又有下界、且每个 $H^n(x)$ 都是 $\mathcal C^{\heartsuit}$ 中紧对象的对象 $x$。
\end{lem}

\begin{defn}{A.6.6}\label{defn:a4-17-A-6-6}
若 Grothendieck 阿贝尔范畴紧生成，且其紧对象所成全子范畴是阿贝尔的，则称它是局部凝聚的（locally coherent）。（关于这一概念的若干等价刻画，例如见 [Her97, Thm.~1.6]。）
\end{defn}

下述定义取自 [BZNP17, Def.~6.2.7]。

\begin{defn}{A.6.7}\label{defn:a4-17-A-6-7}
称容许滤过余极限的稳定 $\infty$-范畴上的 $t$-结构为凝聚的（coherent），是指它右完备、与滤过余极限相容，并且 $\mathcal C^{\heartsuit}$ 局部凝聚。
\end{defn}

由于局部凝聚阿贝尔范畴的紧对象组成阿贝尔子范畴（按定义），凝聚 $t$-结构特别满足引理 A.6.5 的等价条件。

现在我们（仿照 [BZNP17, §6.3]）定义凝聚 $t$-结构何谓正则的（regular）。首先注意：若 $\mathcal C$ 带有凝聚 $t$-结构，则由假设 $\mathcal C$ 容许滤过余极限，存在典范函子
\begin{equation}\tag{A.6.8}\label{eq:a4-17-A-6-8}
\Ind\,\Coh(\mathcal C)\longrightarrow\mathcal C.
\end{equation}
该函子连续（本质上由构造可知）。然后作如下定义。

\begin{defn}{A.6.9}\label{defn:a4-17-A-6-9}
称 $\mathcal C$ 上的凝聚 $t$-结构为正则的，是指 $\mathcal C$ 由其凝聚对象紧生成（意思是：$\mathcal C$ 的凝聚对象恰是 $\mathcal C$ 的紧对象，且它们生成 $\mathcal C$），或者（等价地）函子 (A.6.8) 是等价。
\end{defn}

\begin{rem}{A.6.10}\label{rem:a4-17-A-6-10}
[BZNP17, Defn.~6.3.1] 中正则 $t$-结构的定义并不假定该 $t$-结构凝聚，也没有显然地与定义 A.6.9 一致；不过当该 $t$-结构凝聚时，由 [BZNP17, Prop.~6.3.2] 两个定义一致。
\end{rem}

\ysrc{368}

若 $\mathcal C$ 带有凝聚 $t$-结构，则沿 (A.6.8) 拉回该 $t$-结构，便给 $\Ind\Coh(\mathcal C)$ 赋予一个 $t$-结构；后者仍是凝聚的，而且实际上是正则的 [BZNP17, Prop.~6.3.2]。函子 (A.6.8) 按构造是 $t$-正合的。对任意 $n$，它限制为一个等价
\[
\Ind\Coh(\mathcal C)^{\geq n}\xrightarrow{\sim}\mathcal C^{\geq n}.
\]

若 $\mathcal C$ 带有凝聚 $t$-结构，则其左完备化 $\wh{\mathcal C}$ 上的 $t$-结构仍然是凝聚的 [BZNP17, Prop.~6.3.4]。同样，典范函子 $\mathcal C\to\wh{\mathcal C}$ 诱导等价
\[
\mathcal C^{\geq n}\xrightarrow{\sim}\wh{\mathcal C}^{\geq n}.
\]
于是 $\Ind\Coh(\mathcal C)$ 的构造与 $\wh{\mathcal C}$ 的构造，是「填充」$\mathcal C^+=\bigcup_n\mathcal C^{\geq n}$ 的两种不同方式，填充方向朝左。

\begin{rem}{A.6.11}\label{rem:a4-17-A-6-11}
若 $\mathcal C$ 上的 $t$-结构左分离（例如因为它左完备；见注 A.5.3），则函子 (A.6.8) 的核恰由 $\Ind\Coh(\mathcal C)^{-\infty}$ 组成。在 (A.6.8) 不是等价的例子中（即 $\mathcal C$ 上的 $t$-结构不正则时），常见的情形是 $\Ind\Coh(\mathcal C)^{-\infty}$ 非零，即 $\Ind\Coh(\mathcal C)$ 上的 $t$-结构不左分离。
\end{rem}

为后文引用，我们记下关于 $t$-结构与 Ind-范畴的下述一般事实 [AGH19, Prop.~2.13]。

\begin{prop}{A.6.12}\label{prop:a4-17-A-6-12}
若 $\mathcal C$ 是带有 $t$-结构的稳定 $\infty$-范畴，则 $\Ind\mathcal C$ 继承一个 $t$-结构，其特征是：$\mathcal C\hookrightarrow\Ind\mathcal C$ 为 $t$-正合，且含入 $\mathcal C^{\geq0}\to\mathcal C$ 诱导等价 $\Ind(\mathcal C^{\geq0})\xrightarrow{\sim}\Ind(\mathcal C)^{\geq0}$。若 $\mathcal C$ 上的 $t$-结构还是右有界的，则 $\Ind\mathcal C$ 上的 $t$-结构右完备。
\end{prop}

\subsection*{A.7.\quad 导出函子与伴随}

若 $F:\mathcal A\to\mathcal B$ 是 Grothendieck 阿贝尔范畴之间的加性函子，则 $F$ 诱导相应上链复形范畴之间的函子 $\mathrm{CoCh}(\mathcal A)\to\mathrm{CoCh}(\mathcal B)$，从而诱导函子 $\underline F:\mathrm{CoCh}(\mathcal A)\to D(\mathcal B)$。用 $Q:\mathrm{CoCh}(\mathcal A)\to D(\mathcal A)$ 记局部化函子。按照 [Cis19, §7.5.23]（它把同伦范畴层面上的一个定义「提升」到 $\infty$-范畴的设定；该定义可追溯到 Quillen，见 [Qui67, Def.~1.4.1]，较近的处理见 [Cis19, Def.~2.3.1]），$F$ 的右导出函子是一个函子 $RF:D(\mathcal A)\to D(\mathcal B)$，它配有自然变换 $\underline F\to RF\circ Q$，该自然变换表示函子 $\mathrm{Maps}(F,Q_*(-))$。更确切地说，对任意其他函子 $G:D(\mathcal A)\to D(\mathcal B)$ 连同一个自然变换 $\underline F\to G\circ Q$，存在唯一（在可缩选择空间意义下）的自然变换 $RF\to G$，它给出给定的 $\underline F\to G\circ Q$。

\begin{prop}{A.7.1}\label{prop:a4-17-A-7-1}
设 $F:\mathcal A\to\mathcal B$ 是 Grothendieck 阿贝尔范畴之间的正合函子，且与余极限相容；设 $G:\mathcal B\to\mathcal A$ 记 $F$ 的右伴随（由伴随函子定理存在）。则
\begin{enumerate}[label=(\arabic*)]
\item $F$ 诱导 $t$-正合连续函子 $D(\mathcal A)\to D(\mathcal B)$，我们仍记作 $F$。
\item 右导出函子 $RG:D(\mathcal B)\to D(\mathcal A)$ 存在，并且是 $F$ 的右伴随。
\end{enumerate}
\end{prop}

\ysrc{369}

\begin{proof}
由 $D(\mathcal A)$ 作为上链复形范畴 $\mathrm{CoCh}(\mathcal A)$ 关于拟同构的 $\infty$-范畴局部化这一描述，$F$ 诱导 $t$-正合函子 $D(\mathcal A)\to D(\mathcal B)$ 立得。一旦证明了 $F$ 是左伴随，$F$ 便自动连续。

于是只需证第 \textup{(2)} 部分；我们断言它可由 [Cis19, Thm.~7.5.30] 推出。事实上，$\mathrm{CoCh}(\mathcal A)$（同样 $\mathrm{CoCh}(\mathcal B)$）容许 [Lur17, Prop.~1.3.5.3] 给出的模型结构，而由 [Lur17, Prop.~1.3.5.15]，$D(\mathcal A)$ 是该模型范畴的底层 $\infty$-范畴。现在，因 $F:\mathcal A\to\mathcal B$ 正合，它保持一切拟同构，而其延拓 $F:D(\mathcal A)\to D(\mathcal B)$ 显然就是 $F$ 的左导出函子 $LF$（另见 [Cis19, Lem.~7.5.24]）。由 [Cis19, Thm.~7.5.30]，我们只需证明 $G:\mathrm{CoCh}(\mathcal B)\to\mathrm{CoCh}(\mathcal A)$ 把纤维性上链复形之间的拟同构送到拟同构。而由 [Lur17, Prop.~1.3.5.14]，纤维性上链复形之间的拟同构必为同伦等价，从而被 $G$ 送到同伦等价，特别地送到拟同构，即所求。左伴随，故连续。\tn{底本 p.~369 此处原文作「as required. left adjoint, it is continuous.」。后半句与本节开头「Once we have shown that $F$ is a left adjoint, then $F$ is automatically continuous」重复，显系编辑过程中残留的衍文；译者照录，以存底本之旧。}
\end{proof}

\begin{exam}{A.7.2}\label{exam:a4-17-A-7-2}
命题 A.7.1 的一个典型应用场景是：$\mathcal B$ 是 Grothendieck 范畴，$\mathcal A$ 是局部化子范畴，即 Serre 子范畴，并且对 $\mathcal B$ 中任意直和（从而任意余极限）的构造封闭。此时 $\mathcal A$ 也是 Grothendieck 阿贝尔范畴，我们可以取 $F:\mathcal A\to\mathcal B$ 为含入。

虽然按构造 $F$ 全忠实，但它延拓成函子 $D(\mathcal A)\to D(\mathcal B)$ 之后一般不必全忠实（就我们所知）。不过在这个方向上我们有下述正面结果。（该命题第 \textup{(1)} 条之对偶版本——即对 $D^-$ 而非 $D^+$、涉及投射对象而非内射对象的版本——在三角范畴中的类似物，是 [Sta, Tag 0FCL] 的推论。）
\end{exam}

\begin{prop}{A.7.3}\label{prop:a4-17-A-7-3}
接上文的讨论，设 $F:\mathcal A\hookrightarrow\mathcal B$ 是局部化子范畴到 Grothendieck 阿贝尔范畴的含入，并考虑其 $t$-正合延拓 $D^+(\mathcal A)\to D^+(\mathcal B)$ 与 $D(\mathcal A)\to D(\mathcal B)$。
\begin{enumerate}[label=(\arabic*)]
\item 若 $F$ 满足命题 A.5.7 陈述中的等价性质，则延拓 $D^+(\mathcal A)\to D^+(\mathcal B)$ 仍全忠实，其本质像等于全子范畴 $D^+_{\mathcal A}(\mathcal B)$，后者由上同调落在 $\mathcal A$ 中的对象组成。
\item 假设延拓 $D^+(\mathcal A)\to D^+(\mathcal B)$ 全忠实。再假设 $\mathcal B$ 中积的构造正合，且命题 A.7.1 的导出右伴随 $RG$ 有有限上同调维数。则延拓 $D(\mathcal A)\to D(\mathcal B)$ 仍全忠实，本质像等于 $D_{\mathcal A}(\mathcal B)$（下标 $\mathcal A$ 同样表示由上同调落在 $\mathcal A$ 中的对象组成的全子范畴）。
\end{enumerate}
\end{prop}

\begin{proof}
$D(\mathcal B)$ 上的 $t$-结构诱导 $D^+_{\mathcal A}(\mathcal B)$ 上的 $t$-结构，其中心恰为 $\mathcal A$。若假定 $F$ 满足命题 A.5.7 陈述中的等价性质，则该命题给出 $t$-正合等价
\[
D^+(\mathcal A)\xrightarrow{\sim}D^+_{\mathcal A}(\mathcal B),
\]
它延拓了中心上的恒等函子。该命题的本质唯一性陈述表明，这个等价必本质上与 $F$ 的典范延拓一致。这就证明了 \textup{(1)}。

\ysrc{370}

现在假设 $F$ 到 $D^+(\mathcal A)$ 的延拓全忠实；如上，我们继续把这个延拓以及到 $D(\mathcal A)$ 的延拓都记作 $F$。对 $D(\mathcal A)$ 的任意对象 $X$，伴随的单位给出自然态射 $X\to RG(F(X))$。于是若 $n$ 是任意整数，把这一态射用于正合三角 $\tau^{<n}X\to X\to\tau^{\geq n}X$（并回忆 $F$ 是 $t$-正合的），便得到正合三角之间的态射
\begin{equation*}
\begin{tikzcd}
\tau^{<n}X \arrow[r] \arrow[d] & X \arrow[r] \arrow[d] & \tau^{\geq n}X \arrow[d] \\
RG\bigl(\tau^{<n}F(X)\bigr) \arrow[r] & RG\bigl(F(X)\bigr) \arrow[r] & RG\bigl(\tau^{\geq n}F(X)\bigr)
\end{tikzcd}
\end{equation*}

由于假定函子 $F$ 在 $D^+(\mathcal A)$ 上全忠实，我们看到右边的竖直箭头是同构。当然，$\tau^{<n}X$ 在次数 $\geq n$ 处上同调为零；而若 $a\geq0$ 记 $RG$ 的上同调振幅，则 $RG(\tau^{<n}F(X))$ 在次数 $\geq n+a$ 处上同调为零。于是若 $i\geq n+a$，便得到交换方块
\begin{equation*}
\begin{tikzcd}
H^i(X) \arrow[r, "\sim"] \arrow[d] & H^i(\tau^{\geq n}X) \arrow[d, "\sim"] \\
H^i\bigl(RG(F(X))\bigr) \arrow[r, "\sim"] & H^i\bigl(RG(\tau^{\geq n}F(X))\bigr)
\end{tikzcd}
\end{equation*}
其中（如图中所示）除左边的竖直箭头外，其余箭头都已知是同构。因此这个箭头也是同构。

由于 $n$ 任意，我们发现伴随的单位 $X\to RG(F(X))$ 在所有次数上诱导上同调的同构，从而是 $D(\mathcal A)$ 中的同构。于是 $RG$ 是 $F$ 的左拟逆，故 $F$ 全忠实，即所断言。

于是 $D(\mathcal A)$ 等价于 $D(\mathcal B)$ 的一个全子范畴，后者显然包含在 $D_{\mathcal A}(\mathcal B)$ 中。到这一步，很可能有一个直接论证可以说明 $D_{\mathcal A}(\mathcal B)$ 实际上等于 $F$ 的本质像，但我们用一个稍绕的方式完成论证。

首先注意：由于 $\mathcal B$ 中积的构造正合，范畴 $D(\mathcal B)$ 容许积，且这些积的构造是 $t$-正合的。命题 A.5.4 表明 $D(\mathcal B)$ 还左完备。由此易见 $D_{\mathcal A}(\mathcal B)$ 也左完备。现在，若 $\{X_n\}_{n=0}^{\infty}$ 是 $D(\mathcal A)^{\leq0}$ 中的对象序列，则 $\prod_{n=0}^{\infty}F(X_n)$ 是 $D(\mathcal B)^{\leq0}$ 的对象，从而 $RG\bigl(\prod_{n=0}^{\infty}F(X_n)\bigr)$ 是 $D(\mathcal A)^{\leq a}$ 的对象（其中 $a$ 如上，记 $RG$ 的上同调振幅）。此外，利用 $F$ 与 $RG$ 的伴随性以及 $F$ 的全忠实性可以验证：$RG\bigl(\prod_{n=0}^{\infty}F(X_n)\bigr)$ 是诸 $X_n$ 在 $D(\mathcal A)$ 中的积。于是 $D(\mathcal A)$ 中可数积存在，且这些积的构造具有有界上同调振幅。再次应用命题 A.5.4，便知 $D(\mathcal A)$ 也左完备。


% 本片止于印刷页 370（PDF 页 383）。末段为命题 A.7.3 证明的收尾；
% 下片 ch-a4-18 自印刷页 371 起。
% 命题 A.7.3 的证明跨片：\end{proof} 移至 ch-a4-18.tex 中等价性论证之后。
